The core theorems are finite rewriting statements. Stronger readings of
the same formalism need their own certificates; the table lists them.

\begin{center}
\small
\begin{tabular}{@{}p{0.33\linewidth}p{0.56\linewidth}@{}}
\textbf{Desired statement} & \textbf{Required certificate} \\
\hline
Overlap network is a code & finite constraint-code data of Proposition~\ref{prop:constraint-code} \\
Global repair is a physical normal-form map & \(\operatorname{Rep}_\lambda=\overline{\operatorname{nf}}_\lambda:Q\to Q\), total, idempotent, schedule-independent, landing in \(C_Q\) \\
Repair converges & finite exact descent; confluence additionally needs semantic-complete transactions, coherent canonical aggregate gluing and repair completeness \\
Distance equals min-cut & topological-code certificate with homological logicals and matching systole/min-cut geometry \\
Corrects \(t\) corrupted carriers & code projector, error family, Knill--Laflamme condition and certified \(t<d/2\) distance bound \\
Exponential convergence & declared transfer or channel operator with stationary projection and spectral gap \\
Long-run noisy approximate consensus & fair-block contraction certificate \((\lambda,\varepsilon,A,\beta,L)\); deterministic block endpoints for the calendar-time constant \\
Cross-view BFT safety and liveness & the certificate, lock, quorum, authentication and timing contract of Definition~\ref{def:qbft-style} and (A1)--(A6) \\
\end{tabular}
\end{center}

\subsection*{Byzantine agreement under a prepared-lock contract}

The Byzantine fault model is that of Lamport, Shostak and Pease~\cite{Lamport1982}.

\begin{definition}[QBFT-style protocol]
\label{def:qbft-style}
A consensus protocol is \emph{QBFT-style} in this analysis if it satisfies
the following five structural properties.  Same-view safety uses
(P1)--(P2), cross-view safety additionally uses the lock and monotone-view
parts of (P1) and (P4), and the liveness clause uses (P3)--(P5) together
with the timing and connectivity assumptions below.
\begin{enumerate}[(P1)]
  \item \textbf{One-vote-per-view.}
        Each nonfaulty node supports at most one value per view number.
        If the protocol has prepare, prepared-certificate acceptance, and
        commit phases, its messages in all three phases refer to that same
        value.  A node that has supported a value in
        view~$v$ ignores any later conflicting request in view~$v$.
        View participation is monotone: after entering view \(w\), a
        nonfaulty node never signs a prepare, prepared-certificate acceptance,
        or commit message for \(v<w\), and never adopts a durable lower-view
        lock.
  \item \textbf{Certificate semantics.}
        A raw prepared certificate
        \(\operatorname{PC}(v,x)\) contains \(q\) distinct, unforgeable,
        authenticated prepare votes for value \(x\) in view \(v\).
        A lock certificate \(\operatorname{LC}(v,x)\) contains that
        \(\operatorname{PC}(v,x)\) together with \(q\) distinct authenticated
        acceptance acknowledgements. These acknowledgements are provisional:
        they do not change the acceptor's voting lock. A raw prepared
        certificate without \(q\) acceptance acknowledgements cannot justify a
        new view or a commit vote. A nonfaulty validator that receives the
        assembled \(\operatorname{LC}(v,x)\) verifies it and records the
        durable lock \((v,x)\) before signing a commit vote. A decision
        certificate contains \(q\) distinct authenticated commit votes, each
        referring to the same valid \(\operatorname{LC}(v,x)\).  In the
        classical exact-size case \(q=2f+1\).  For larger validator sets at
        fixed fault budget, the threshold must scale so that (A6) holds.
        A nonfaulty observer finalises a value only after accepting such a
        valid decision certificate.
  \item \textbf{DLS-style view-change.}
        If no certificate is produced within a timeout, every nonfaulty node
        increments the view number by one and a new leader is selected by a
        fixed fair deterministic rule.  Once the post-\textrm{GST} timeout
        bound of (A1) is active, the pacemaker brings all nonfaulty nodes
        through successive views and reaches a nonfaulty leader within at
        most \(f+1\) consecutive views.
  \item \textbf{Prepared-certificate dissemination, commit lock, and
        justified new view.}
        The assembler disseminates a valid \(\operatorname{PC}(v,x)\) and
        obtains the \(q\) authenticated acceptance acknowledgements of
        \(\operatorname{LC}(v,x)\), then disseminates the assembled lock
        certificate before collecting commits.
        Each validator's authenticated view-change message carries its highest valid
        lock certificate, if any.  A new leader collects \(q\) such
        messages.  If it issues a proposal, that proposal carries the value
        of their highest-view lock certificate; it may carry a fresh
        value only if none is present.  The proposal carries the complete
        new-view justification.
        Every nonfaulty validator verifies the signatures, quorum size, and
        highest-certificate selection.  An unlocked validator may support a
        fresh proposal only when the verified justification contains no
        lock certificate.  Otherwise it votes only for its current lock
        value, except that the value selected by a valid \(q\)-message new-view
        justification supersedes its local assembled-lock-certificate lock.
        This authorizes the prepare vote without adopting a lower-view durable
        lock; receipt of the assembled current-view lock certificate records
        the replacement before commit. This rule applies
        without a local test for whether an earlier commit set reached \(q\):
        it reconciles every predecision lock, while the proof below shows that
        quorum intersection makes the selected value compatible whenever a
        decision certificate exists.
  \item \textbf{Post-\textrm{GST} nonfaulty progress.}
        In a sufficiently long post-\textrm{GST} view led by a nonfaulty
        node, that leader collects the required new-view messages and emits
        the proposal selected by (P4).  Every nonfaulty validator that
        receives a valid proposal promptly sends its prepare vote.  The
        nonfaulty leader disseminates the resulting valid prepared
        certificate; every nonfaulty validator verifies it and returns a
        provisional acceptance acknowledgement. After assembling the lock
        certificate, the leader disseminates it; every nonfaulty recipient
        verifies it, records or advances its durable lock, and sends its
        commit vote.
        After accepting a valid decision certificate it
        finalises and relays that certificate.  Each of these fixed protocol
        phases completes within \(O(\Theta)\), with \(\Theta\) defined in
        (A1).
\end{enumerate}
The Istanbul BFT / QBFT protocol family \cite{Moniz2020,Saltini} motivates
this pattern; (P1)--(P5), including their exact certificate, pacemaker, lock,
and progress semantics, must be checked against an implementation rather
than inferred from the family name.
\end{definition}

\paragraph{Assumptions A1--A6.}
\begin{enumerate}[(A1)]
  \item \textbf{Partial synchrony~\cite{DLS88}.}
        Fixed but initially unknown bounds $\Delta$ (per-edge message delay)
        and $\Phi$ (local processing time) hold after \textrm{GST}. Together
        with the certified routes of (A4), they supply an effective
        end-to-end protocol-phase bound \(\Theta\) that includes route length,
        processing, and the fixed number of message steps in one phase.
        Pacemaker timeouts eventually exceed the corresponding per-view
        bound.
        \emph{Safety} holds without extra timing assumptions;
        \emph{liveness} holds after the Global Stabilisation Time (\textrm{GST}).
  \item \textbf{Byzantine fault model.}
        At most $f$ observers behave arbitrarily; the remaining $n-f$ are nonfaulty.
  \item \textbf{Classical exact sizing for this fixed-quorum theorem.}
        $n = 3f+1$, so that $q=2f+1$ certificates have a nonfaulty overlap
        witness.  The usual resilience condition $n\ge 3f+1$ is necessary for
        the fault model, but when $n>3f+1$ a fixed $q=2f+1$ quorum is
        insufficient for the overlap used in the safety proof; one must either
        impose (A6) directly or choose a threshold $q$ with
        $2q\ge n+f+1$.
  \item \textbf{Strong quorum connectivity.}
        Every quorum $Q$ with $|Q|=q$ is strongly connected within~$G$:
        for any $u,v\in Q$ there is a directed path in~$G$ contained
        entirely in~$Q$.
        This is strictly stronger than requiring the overlap graph of
        quorums to be connected, and is needed to propagate signed votes
        within a quorum.
  \item \textbf{Message authentication.}
        All messages carry unforgeable digital signatures.
  \item \textbf{OPH quorum overlap.}
        Any two quorums $Q_a, Q_b$ of size $q$ satisfy
        $|Q_a\cap Q_b|\geq f+1$.  For $q=2f+1$ this is guaranteed by the
        exact sizing $n=3f+1$ in (A3); for general $n$ it is the separate
        threshold condition $2q\ge n+f+1$, not a consequence of
        $n\ge3f+1$ alone.  The liveness clause additionally requires
        \(q\le n-f\), so a certificate can be formed without Byzantine
        participation.
\end{enumerate}

\begin{theorem}[QBFT safety and liveness]
\label{thm:qbft-safety}
Under (A1), (A2), and (A4)--(A6), any consensus protocol satisfying
(P1)--(P5) of Definition~\ref{def:qbft-style} and run over the
\textrm{OPH} observer graph satisfies:
\begin{enumerate}[(i)]
  \item \textbf{Cross-view safety.} No two nonfaulty observers finalise
        conflicting patch states, whether in the same view or in different
        views.
  \item \textbf{Liveness.} From the first post-\textrm{GST} view in which the
        pacemaker timeout bound is active, every nonfaulty observer finalises
        within $O((f+1)\Theta)$ wall-clock time, provided \(q\le n-f\).
        No bound from \textrm{GST} to timeout activation is claimed.
  \item \textbf{Certificate-contract threshold.} Joint feasibility of the
        safety threshold \(2q\ge n+f+1\) and Byzantine-independent liveness
        threshold \(q\le n-f\) forces \(n\ge3f+1\); the exact-size choice
        \(n=3f+1,q=2f+1\) attains both.
\end{enumerate}
\end{theorem}

\paragraph{Sizing.}
Assumption (A3) is the classical exact-size way to discharge (A6): when
\(n=3f+1\) and \(q=2f+1\) the required overlap is automatic, and for general
\(n\) (A6) must be checked directly. The prepared-lock rule proves safety
only for the protocol just defined, never from the name ``QBFT.'' Under
these assumptions no two nonfaulty observers finalise conflicting patch
states.

\begin{proof}
\textit{Same-view safety.}
Suppose $O_a$ and $O_b$ finalise $s_a\neq s_b$ in one view.
By (P2), each required a certificate of $q$ votes: sets $Q_a,Q_b$.
By (A6), $|Q_a\cap Q_b|\geq f+1$.  In the classical exact-size case this is
the inclusion-exclusion calculation
$|Q_a\cap Q_b|\geq(2f+1)+(2f+1)-(3f+1)=f+1$.
By (A2), at most $f$ are Byzantine, so $Q_a\cap Q_b$ contains a
nonfaulty $O^*$.
By (P1), $O^*$ voted for at most one value. Contradiction.
The identical overlap argument applies to two prepared certificates in the
same view, using their prepare-vote quorums.

\textit{Cross-view safety.}
Suppose \(\operatorname{DC}(v,x)\) and
\(\operatorname{DC}(w,y)\) are conflicting decision certificates with
\(v<w\). The first certificate refers to a valid
\(\operatorname{LC}(v,x)\). Let \(C_x\) be its \(q\) committers. The
nonfaulty subset \(L_x\subseteq C_x\) has size at least \(q-f\), and every
member of \(L_x\) received the assembled lock certificate and recorded
\((v,x)\) before committing. Every size-\(q\) prepare or view-change quorum \(P\)
intersects that actual lock set because
\[
 |L_x\cap P|
 \ge (q-f)+q-n
 =2q-n-f
 \ge1
\]
by (A6).

Assume for contradiction that a raw prepared certificate for a value
different from \(x\) exists above view \(v\), and choose one with least view
\(r>v\). Its proposal carries a valid \(q\)-message new-view justification
\(J_r\). Since \(|L_x|\ge q-f\) and \(|J_r|=q\), choose a nonfaulty
\(O^*\in L_x\cap J_r\). Validator \(O^*\) signed a commit for
\(\operatorname{LC}(v,x)\), and later sent its authenticated view-\(r\)
change. Monotone view participation in (P1) orders its view-\(v\) lock and
commit before that view-\(r\) message: it cannot send a view-\(r\) change and
subsequently return to commit in view \(v\). Its view-change message therefore
reports \(\operatorname{LC}(v,x)\), or a later valid lock certificate. By
minimality of \(r\), every raw prepared certificate underlying such a later
reported lock certificate below view \(r\) carries \(x\). Hence \(O^*\)'s
report carries \(x\).

The justification \(J_r\) therefore cannot be certificate-free. It also
cannot contain a conflicting highest lock certificate, because that
certificate contains a conflicting raw prepared certificate in a view below
\(r\), contradicting the minimality of \(r\). Rule (P4) consequently selects
\(x\), so nonfaulty validators do not prepare the alleged conflicting value
in view \(r\).
Thus no higher conflicting prepared certificate exists, and (P2) excludes
the alleged \(\operatorname{DC}(w,y)\).

This induction is also the lock-transfer invariant for view change. Every
valid \(q\)-message new-view justification after a decision contains a
nonfaulty report from the decision's \(q-f\) commit-lock set, and its highest
valid report carries the decided value.

\textit{Liveness and certificate threshold.}
From the first post-\textrm{GST} view whose pacemaker timeout is active,
(A1) makes every failed-view timeout and every fixed protocol phase
\(O(\Theta)\). Rule (P3) reaches a nonfaulty leader within at most \(f+1\)
consecutive views. Since \(q\le n-f\), nonfaulty validators can
supply all \(q\) new-view, prepare, prepared-certificate acceptance, and
commit messages without Byzantine participation.

Consider the valid \(q\)-message justification used by the successful new
view. If it contains no lock certificate, (P4) selects a fresh value.
Otherwise it selects its unique highest-view lock certificate; same-view
uniqueness follows from the prepare-quorum overlap. Every nonfaulty validator
may replace a local assembled-lock-certificate predecision lock with that
valid selection, including a
lock whose local commit set never reached a decision certificate. If a
decision certificate exists, the \(q-f\) nonfaulty commit-lock set proved
above intersects the \(q\)-message justification. Its report forces the
selected lock certificate to carry the decided value, since the safety
induction excludes a conflicting certificate. Thus at least \(q\) nonfaulty
validators can support the selected proposal in either case, even when some
of their orphan locks were omitted from the new-view messages.

Rule (P5) executes the proposal, prepare, provisional acceptance,
lock-certificate dissemination, commit, and decision-relay phases, so every
nonfaulty observer finalises within \(O((f+1)\Theta)\) from that activation
point. Eventual timeout activation yields eventual liveness, but the theorem
does not bound its delay after \textrm{GST}.

Finally, \(q\le n-f\) and \(2q\ge n+f+1\) imply
\[
 2(n-f)\ge n+f+1,
\]
hence \(n\ge3f+1\) (with the integer convention specified in (A6)).
At \(n=3f+1\), \(q=2f+1\) attains both inequalities. This is a threshold
statement for the declared certificate contract, not a universal resilience
claim for every authenticated protocol model.

\textit{Note on FLP.}
Fischer, Lynch, Paterson~\cite{FLP85} is an impossibility result for
fully asynchronous systems; it does not bear on achievability under
partial synchrony (A1).
\end{proof}

\paragraph{Executable receipt and negative controls.}
A finite executable check exhausts all certificate and new-view quorum pairs
for \((n,f,q)=(1,0,1),(4,1,3),(6,1,4),(7,2,5)\); the third checks the
general threshold with \(q<n-f\). It checks the certificate-overlap witness,
independent prepare-signer, provisional-acceptor and committer quorums, the
\(q-f\) lower bound on nonfaulty commit locks, their transfer into every
new-view quorum, and the next-view lock transition against every candidate
prepare quorum, including an omitted higher orphan lock, together with
post-GST traces for every nonfaulty reference leader. Separate finite traces
remove each load-bearing safety or liveness premise. The lock-free
\(n=4,f=1,q=3\) trace decides \(x\) from \(\{B,N_1,N_2\}\) in view zero and
\(y\ne x\) from \(\{B,N_2,N_3\}\) in view one while respecting
one-value-per-validator in each view. The receipt is not a general
protocol-model checker; the arbitrary-view induction is the proof above.

\subsection*{The Petz repair map}

\begin{definition}[OPH Repair Map: Petz form]
\label{def:oph-repair-map}
Let $\sigma\in\mathcal{D}(\mathcal{H})$ be a full-rank reference state and
$\mathcal{N}:\mathcal{B}(\mathcal{H})\to\mathcal{B}(\mathcal{K})$ a quantum channel.
The \emph{OPH repair map} is
\[
  \mathcal{R}_{\sigma,\mathcal{N}}(\rho)
  := \sigma^{1/2}\,
     \mathcal{N}^\dagger\!\bigl(
       \mathcal{N}(\sigma)^{-1/2}\,\rho\,\mathcal{N}(\sigma)^{-1/2}
     \bigr)
     \,\sigma^{1/2},
\]
where $\mathcal{N}^\dagger$ is the adjoint channel and inverses are
taken on $\mathrm{supp}(\mathcal{N}(\sigma))$.
\end{definition}

The closest-point trace-distance projection
\(\arg\min_{\tau\in\mathcal S}\tfrac12\|\rho-\tau\|_1\) is a different object:
it is variational and not CPTP in general. All properties below refer to
Definition~\ref{def:oph-repair-map}.

\begin{proposition}[Petz map CPTP on its domain]
\label{prop:petz-cptp}
Let $\sigma$ have full support on $\mathcal{H}$.
\begin{enumerate}[(a)]
  \item $\mathcal{R}_{\sigma,\mathcal{N}}$ is completely positive.
  \item $\mathcal{R}_{\sigma,\mathcal{N}}$ is trace-preserving on
        $\mathrm{supp}(\mathcal{N}(\sigma))$, i.e., on inputs $\rho$
        for which $\mathcal{N}(\sigma)^{-1/2}\rho\,\mathcal{N}(\sigma)^{-1/2}$
        is well-defined.
  \item If additionally $\mathcal{N}(\sigma)$ has full rank on $\mathcal{K}$,
        then $\mathcal{R}_{\sigma,\mathcal{N}}$ is \textrm{CPTP} on all
        of $\mathcal{B}(\mathcal{K})$.
\end{enumerate}
If $\mathcal{N}(\sigma)$ is \emph{not} full rank on $\mathcal{K}$,
then either (i)~the domain must be restricted to
$\mathrm{supp}(\mathcal{N}(\sigma))$, or
(ii)~pseudoinverses must replace the inverses (generalised Petz map;
cf.~\cite{Junge2018}), or
(iii)~a regularisation
$\mathcal{N}(\sigma)\mapsto\mathcal{N}(\sigma)+\varepsilon\mathbf{1}$
must be introduced.
Note that full-rank $\sigma$ does \emph{not} prevent $\mathcal{N}(\sigma)$
from being rank-deficient: the channel may map the support of $\sigma$
into a strict subspace of $\mathcal{K}$.
In the \textrm{OPH} setting, whether $\mathcal{N}(\sigma)$ is full rank
depends on the specific overlap channel and must be verified for the chosen analytic channel
model. The finite OPH repair theorem instead uses the declared Lyapunov descent law on a finite
patch net.
\end{proposition}

\begin{proof}
Complete positivity follows from composing three \textrm{CP} operations:
\begin{enumerate}[(i)]
  \item sandwiching by
        $\mathcal{N}(\sigma)^{-1/2}(\cdot)\mathcal{N}(\sigma)^{-1/2}$
        on $\mathrm{supp}(\mathcal{N}(\sigma))$;
  \item $\mathcal{N}^\dagger$;
  \item sandwiching by $\sigma^{1/2}(\cdot)\sigma^{1/2}$.
\end{enumerate}
Trace preservation in the full-rank case:
Petz~\cite{Petz1986}; Fagnola--Umanit\`{a}~\cite{FagnolaUmanita2010}.
\end{proof}

Analytic refinements of the repair map are separate conditions. A strictly
contractive quotient repair map \(T\) with coefficient \(\lambda<1\) has at
most one fixed point; when one exists, the iterates of an implemented map
within \(\varepsilon\) of \(T\) approach it to within
\(\varepsilon/(1-\lambda)\). Exponential
convergence of a declared transfer or channel operator requires a spectral
gap \(\|\mathcal T^t-\Pi_\star\|\le Ce^{-\delta t}\) on the nonstationary
subspace. Neither condition is used by the finite normal-form theorem, where
termination follows from strict descent of accepted repairs. Strong fairness
imposes no probability space on schedules, no spectral gap and no
message-delay bound, so it yields no probability-one, exponential or
wall-clock statement~\cite{FLP85}.

\subsection*{Codes}

\begin{theorem}[No free min-cut for bare overlap graphs]
\label{thm:no-free-mincut}
Let \(G=(V,E)\) be any connected graph with \(|V|\ge2\). The graph \(G\) alone does not determine
the Hamming distance of the consistency set \(C\) of a finite overlap net on \(G\). In particular,
the same graph can realize a binary constraint code of distance \(1\) or a binary repetition code
of distance \(|V|\).
\end{theorem}

\begin{proof}
Set \(S_i=\{0,1\}\) for every vertex. For the repetition realization, choose \(I_e=\{0,1\}\) and
let both endpoint readouts be the identity. Then every edge imposes \(x_i=x_j\), so the only
global codewords are \(00\cdots0\) and \(11\cdots1\), whose Hamming distance is \(|V|\).

For the trivial-overlap realization, keep the same vertex state spaces but let every endpoint
readout be the constant map to \(0\). Then every binary assignment is globally consistent, so the
minimum Hamming distance among distinct codewords is \(1\). The graph is unchanged. Therefore
distance is a property of the code realization (state spaces, readout maps, logical dictionary,
metric, and error model), not of the bare overlap graph.
\end{proof}

\begin{definition}[Topological-code realization certificate]\label{def:tcert}
An OPH overlap network may be treated as a QECC/topological code only after supplying a tuple
\[
\mathsf{TCert}=
(K,\mathcal H_{\mathrm{phys}},\mathcal H_{\mathrm{code}},
\partial_2,\partial_1,S_X,S_Z,\mathcal L_X,\mathcal L_Z,\mathcal E,\mathcal R),
\]
where \(C_2\xrightarrow{\partial_2}C_1\xrightarrow{\partial_1}C_0\) is a chain complex over
\(\mathbb F_2\), the physical carriers live in \(\mathcal H_{\mathrm{phys}}\), the protected
subspace is \(\mathcal H_{\mathrm{code}}\), \(S_X,S_Z\) are stabilizer or gauge checks,
\(\mathcal L_X,\mathcal L_Z\) are logical-operator classes, \(\mathcal E\) is a declared error
family, and \(\mathcal R\) is a recovery map or recovery family.
\end{definition}

\begin{theorem}[Certified topological-code distance and min-cut]
\label{thm:tcert-distance}
Suppose a certificate \(\mathsf{TCert}\) of Definition~\ref{def:tcert} is supplied, with logical
classes identified as
\[
\mathcal L_X\simeq H_1(K;\mathbb F_2),
\qquad
\mathcal L_Z\simeq H^1(K;\mathbb F_2),
\]
and with boundary conditions excluding lower-weight trivial representatives. Then the certified
distance is
\[
d=
\min\left\{
\min_{\ell\in\mathcal L_X\setminus0}|\ell|,
\min_{\ell^\star\in\mathcal L_Z\setminus0}|\ell^\star|
\right\}.
\]
Only in geometries where this homological systole equals the relevant graph min-cut may one write
\(d=\mathrm{mincut}(G_{\mathrm{OPH}})\) \cite{Kitaev2003,Dennis2002}.
\end{theorem}

\begin{proof}
This is the standard stabilizer/topological-code distance statement once the chain complex,
checks, logical representatives, and boundary conditions are declared. The min-cut equality is an
additional geometric identification of that homological minimum with a graph cut. By
Theorem~\ref{thm:no-free-mincut}, it cannot be inferred from the bare graph.
\end{proof}

\begin{theorem}[QECC resilience under a supplied code certificate]
\label{thm:qecc}
Assume a genuine code subspace \(\mathcal H_{\mathrm{code}}\subseteq\mathcal H_{\mathrm{phys}}\)
with projector \(\Pi\), and let \(\mathcal E_t\) be the declared set of errors supported on fewer
than \(t\) corrupted patches or physical carriers. If
\[
\Pi E_a^\dagger E_b \Pi=\alpha_{ab}\Pi
\qquad
\forall E_a,E_b\in\mathcal E_t,
\]
then there exists a recovery channel correcting all errors in \(\mathcal E_t\)
\cite{KL97}. If the supplied certificate also gives distance \(d\), then all errors of weight
\(t<d/2\) are correctable. No such resilience statement follows from the bare overlap graph.
\end{theorem}
